\documentclass[11pt]{article}

\usepackage[a4paper,margin=1in]{geometry}
\usepackage[T1]{fontenc}
\usepackage[utf8]{inputenc}
\usepackage{lmodern}
\usepackage{microtype}
\usepackage{amsmath,amssymb,amsfonts,mathtools,bm}
\usepackage{booktabs}
\usepackage{array}
\usepackage{enumitem}
\usepackage{xcolor}
\usepackage[most]{tcolorbox}
\usepackage{hyperref}
\usepackage[nameinlink,capitalise]{cleveref}

\hypersetup{
  colorlinks=true,
  linkcolor=blue!55!black,
  citecolor=blue!55!black,
  urlcolor=blue!55!black,
  pdftitle={Bayesian extrapolation of an incrementally changed graph signal with optimal transport regularization},
  pdfauthor={OpenAI}
}

\numberwithin{equation}{section}
\allowdisplaybreaks

\newcommand{\R}{\mathbb{R}}
\newcommand{\N}{\mathbb{N}}
\newcommand{\E}{\mathbb{E}}
\newcommand{\Var}{\operatorname{Var}}
\newcommand{\Cov}{\operatorname{Cov}}
\newcommand{\KL}{\operatorname{KL}}
\newcommand{\diag}{\operatorname{diag}}
\newcommand{\tr}{\operatorname{tr}}
\newcommand{\rank}{\operatorname{rank}}
\newcommand{\argmin}{\operatorname*{arg\,min}}
\newcommand{\argmax}{\operatorname*{arg\,max}}
\newcommand{\ones}{\mathbf{1}}
\newcommand{\eps}{\varepsilon}
\newcommand{\dd}{\mathrm{d}}
\newcommand{\transpose}{^{\mathsf T}}
\newcommand{\pinv}{^{\dagger}}
\newcommand{\given}{\,|\,}
\newcommand{\norm}[1]{\left\lVert #1\right\rVert}
\newcommand{\inner}[2]{\left\langle #1,#2\right\rangle}
\newcommand{\half}{\frac{1}{2}}
\newcommand{\Prob}{\mathbb{P}}

\newtcolorbox{runningexample}[1][]{
  breakable,
  colback=blue!2!white,
  colframe=blue!55!black,
  title=#1,
  fonttitle=\bfseries,
  left=1mm,right=1mm,top=1mm,bottom=1mm
}

\newtcolorbox{heuristicbox}[1][]{
  breakable,
  colback=green!2!white,
  colframe=green!45!black,
  title=#1,
  fonttitle=\bfseries,
  left=1mm,right=1mm,top=1mm,bottom=1mm
}

\newtcolorbox{warningbox}[1][]{
  breakable,
  colback=red!2!white,
  colframe=red!50!black,
  title=#1,
  fonttitle=\bfseries,
  left=1mm,right=1mm,top=1mm,bottom=1mm
}

\title{Bayesian extrapolation of an incrementally changed graph signal\\
\large with graph optimal-transport regularization, derivations, and a running numerical example}
\author{Technical note}
\date{\today}

\begin{document}
\maketitle

\begin{abstract}
We consider $N$ normalized items living in the same feature space.  Each item has a scalar value $x_i$, the inter-dependencies are encoded by a sparse directed weighted graph, the time-$0$ scalar field is known with finite Bayesian precision, and at time $1$ only $n\ll N$ nodes are observed.  The task is to infer both the likely time-$1$ values of the unobserved nodes and their posterior precision levels.  The central modeling choice is to regularize the \emph{increment}
\[
  z=x^1-x^0,
\]
not the absolute time-$1$ field.  The statistically clean default is a Gaussian Bayesian inverse problem on the graph.  A graph optimal-transport (OT) regularizer enters through the local, small-increment, Benamou-Brenier tangent metric: an increment is plausible when it can be generated by a low-energy flux over cheap graph edges, possibly with controlled source/sink creation.  This choice keeps inference linear-Gaussian, so the posterior mean and marginal posterior precisions have explicit formulas.  A nonlinear unbalanced OT model is also given for cases where the scalar values are genuine nonnegative masses.  The exposition is written so that an advanced undergraduate with linear algebra, multivariate Gaussians, and basic optimization can reproduce the derivations.
\end{abstract}

\tableofcontents

\section{The problem in one sentence}

You know a baseline graph signal $\bar x^0\in\R^N$ imperfectly, you observe a small subset of the next signal $x^1\in\R^N$, and you want to fill in the unobserved coordinates while honestly reporting uncertainty.  The graph tells you which increments should influence which other increments.

Let
\[
  V=\{1,\ldots,N\}
\]
be the node set.  The directed weights are $w_{ij}\ge 0$.  A large $w_{ij}$ means that node $j$ is relevant when predicting node $i$, but the matrix need not be symmetric.  The baseline measurement model is
\begin{equation}
  \bar x_i^0=x_i^{0,\star}+\xi_i^0,
  \qquad
  \xi_i^0\sim\mathcal N(0,(p_i^0)^{-1}),
  \label{eq:baseline-measurement}
\end{equation}
where $x^{0,\star}$ is the unknown true baseline and $p_i^0$ is the baseline precision.  At time $1$ we observe only nodes
\[
  S=\{s_1,\ldots,s_n\}\subset V,
  \qquad n\ll N.
\]
The time-$1$ observations are
\begin{equation}
  y_a=x^1_{s_a}+\epsilon_a,
  \qquad
  \epsilon_a\sim\mathcal N(0,r_a^{-1}),
  \qquad a=1,\ldots,n.
  \label{eq:time-one-observation}
\end{equation}
Let $H\in\R^{n\times N}$ be the selection matrix satisfying
\[
  Hx=(x_{s_1},\ldots,x_{s_n})\transpose,
\]
and set
\begin{equation}
  P_0=\diag(p_1^0,\ldots,p_N^0),
  \qquad
  R=\diag(r_1,\ldots,r_n).
  \label{eq:P0-R}
\end{equation}
The output we want is
\begin{equation}
  \widehat x_i^1=\E[x_i^1\given y,\bar x^0],
  \qquad
  \pi_i^+=\frac{1}{\Var(x_i^1\given y,\bar x^0)},
  \qquad i=1,\ldots,N.
  \label{eq:desired-output}
\end{equation}
The scalar $\pi_i^+$ is the \emph{marginal} posterior precision.  It is the reciprocal of the posterior variance of node $i$ alone.

\begin{warningbox}[A frequent precision mistake]
If $Q^+$ is the full posterior precision matrix, then $Q^+_{ii}$ is a \emph{conditional} precision: it measures how tightly $x_i$ is determined if all other coordinates are already known.  The precision requested in \eqref{eq:desired-output} is usually
\[
  \pi_i^+=\frac{1}{K^+_{ii}},
  \qquad K^+=(Q^+)^{-1},
\]
not $Q^+_{ii}$.  In correlated systems $Q^+_{ii}$ is often too optimistic as a marginal uncertainty report.
\end{warningbox}

\begin{runningexample}[Running numerical example: raw data]
We use a six-node graph throughout the note.  The row-normalized transition weights $K_{ij}=w_{ij}/\sum_k w_{ik}$ are
\[
\begin{array}{c|l}
\toprule
\text{node }i & \text{nonzero outgoing probabilities }K_{ij}\\
\midrule
1 & K_{12}=1.000\\
2 & K_{21}=0.333,\quad K_{23}=0.571,\quad K_{25}=0.095\\
3 & K_{32}=0.529,\quad K_{34}=0.471\\
4 & K_{43}=0.375,\quad K_{45}=0.625\\
5 & K_{52}=0.048,\quad K_{54}=0.429,\quad K_{56}=0.524\\
6 & K_{65}=1.000\\
\bottomrule
\end{array}
\]
The baseline values and baseline precisions are
\[
\begin{array}{c|cccccc}
\toprule
 i&1&2&3&4&5&6\\
\midrule
 \bar x_i^0&2.00&2.10&1.95&1.05&0.90&1.10\\
 p_i^0&25&25&16&16&20&20\\
\bottomrule
\end{array}
\]
Thus the baseline standard deviations are $1/\sqrt{p_i^0}$, for example $0.20$ at nodes $1$ and $2$.  At time $1$ we observe
\[
  y_2=2.45,\qquad r_2=100,
  \qquad
  y_5=0.65,
  \qquad r_5=64.
\]
So the observed innovations are
\[
  d_2=y_2-\bar x_2^0=0.35,
  \qquad
  d_5=y_5-\bar x_5^0=-0.25.
\]
The example is deliberately small enough that every matrix can be written down and checked by hand, but it has the same structure as the large sparse case.
\end{runningexample}

\section{Why the increment, not the absolute level, should be regularized}
\label{sec:increment-not-level}

The phrase ``time $1$ is only incrementally different from time $0$'' has a precise modeling consequence.  Write
\begin{equation}
  x^1=x^{0,\star}+m_\Delta+z,
  \label{eq:increment-decomposition}
\end{equation}
where $m_\Delta\in\R^N$ is an optional drift and $z$ is a residual increment.  When there is no learned trend, take $m_\Delta=0$.

A regularizer on $x^1$ itself says that the new absolute field should be graph-smooth, or graph-transportable, or small.  That is often wrong.  The baseline field $x^0$ may contain true sharp features: different clusters may have different levels, and those differences should be preserved.  A regularizer on $z$ says instead:

\begin{quote}
The baseline is the default prediction.  The time-$1$ data may force changes, but the changes should be small and coherent with the graph.
\end{quote}

This distinction is essential.  If node $i$ is unobserved and far from all observed time-$1$ innovations, then the posterior mean should remain close to $\bar x_i^0$.  Its uncertainty should nevertheless reflect the baseline precision $p_i^0$, the expected size of increments, and its graph connection to the observed subset.

\begin{heuristicbox}[Mental picture]
Think of $\bar x^0$ as a terrain map and $z$ as a weather update.  You do not want to smooth the mountains away.  You want to infer where the weather update has moved, diffused, or been created or destroyed.
\end{heuristicbox}

\section{The predictive prior at time $1$}
\label{sec:predictive-prior}

We first build a prior distribution for $x^1$ before conditioning on the new observations.  Suppose the residual increment satisfies
\begin{equation}
  z\sim\mathcal N(0,K_\Delta),
  \qquad
  Q_\Delta:=K_\Delta^{-1},
  \label{eq:increment-Gaussian}
\end{equation}
where $Q_\Delta$ is a graph-informed precision matrix designed later.  Under a locally flat prior for the unknown baseline $x^{0,\star}$, the baseline measurement model \eqref{eq:baseline-measurement} gives
\begin{equation}
  x^{0,\star}\given \bar x^0\sim\mathcal N(\bar x^0,P_0^{-1}).
  \label{eq:baseline-posterior-flat}
\end{equation}
Assume $z$ is independent of the baseline measurement error.  Then from \eqref{eq:increment-decomposition},
\begin{equation}
  x^1\given\bar x^0
  \sim \mathcal N(\mu^-,K^-),
  \qquad
  \mu^-:=\bar x^0+m_\Delta,
  \qquad
  K^-:=P_0^{-1}+K_\Delta.
  \label{eq:predictive-prior}
\end{equation}
The finite precision of the baseline therefore adds variance to the time-$1$ predictive distribution.  If the baseline is effectively exact, replace $P_0^{-1}$ by the zero matrix.

\subsection{Derivation of \texorpdfstring{$K^-=P_0^{-1}+K_\Delta$}{K-=P0^-1+KDelta}}

For clarity, set $m_\Delta=0$.  From \eqref{eq:increment-decomposition},
\[
  x^1=x^{0,\star}+z.
\]
Using \eqref{eq:baseline-posterior-flat}, write
\[
  x^{0,\star}=\bar x^0+u,
  \qquad
  u\sim\mathcal N(0,P_0^{-1}),
\]
where $u$ is independent of $z$.  Therefore
\[
  x^1=\bar x^0+u+z.
\]
The mean is $\bar x^0$, and the covariance is
\[
  \Cov(u+z)=\Cov(u)+\Cov(z)=P_0^{-1}+K_\Delta.
\]
This is \eqref{eq:predictive-prior}.

\begin{warningbox}[Why baseline uncertainty matters]
Ignoring $P_0^{-1}$ makes the posterior overconfident.  In the extreme, if node $i$ is unobserved at time $1$ and disconnected from all observed nodes, the posterior should still remember that the baseline value was measured with finite precision.
\end{warningbox}

\section{Turning a directed graph into valid Bayesian operators}
\label{sec:directed-graph-operators}

A Gaussian precision matrix must be symmetric positive semidefinite, and preferably positive definite after adding local temporal precision.  A raw directed graph Laplacian is generally nonsymmetric, so it is not directly a valid Gaussian precision.  We separate the directed graph into two safe objects:

\begin{enumerate}[leftmargin=2em]
  \item a reversible undirected conductance graph, used for OT geometry and ordinary graph Laplacians;
  \item an optional symmetric directed-residual penalty that still remembers the directionality.
\end{enumerate}

\subsection{Row normalization and stationary mass}

Let
\begin{equation}
  K_{ij}=\frac{w_{ij}}{\sum_k w_{ik}}
  \label{eq:row-normalization}
\end{equation}
when node $i$ has positive out-degree.  If a node has no outgoing edges, either add a small self-loop, add a small teleportation probability, or work on connected components.  If $K$ is irreducible, let $\pi\in\R^N$ be its stationary distribution:
\begin{equation}
  \pi_i>0,
  \qquad
  \sum_i\pi_i=1,
  \qquad
  \pi\transpose K=\pi\transpose.
  \label{eq:stationary-pi}
\end{equation}
The value $\pi_i$ is the long-run mass of node $i$ under the random walk.

\subsection{Reversibilized conductances}

For unordered pairs $\{i,j\}$ define
\begin{equation}
  c_{ij}=\half(\pi_iK_{ij}+\pi_jK_{ji}),
  \qquad i<j.
  \label{eq:reversible-conductance}
\end{equation}
This is the average stationary traffic across the directed edge pair.  It is symmetric by construction.  It is zero if neither direction is present, and it remains positive if only one direction is present.

Let $E$ be the undirected support $\{\{i,j\}:c_{ij}>0\}$.  Choose any orientation of each undirected edge.  The incidence matrix $B\in\R^{N\times |E|}$ has one $+1$ and one $-1$ in each column.  With $C=\diag(c_e)$, the reversible graph Laplacian is
\begin{equation}
  L_{\rm rev}=BCB\transpose.
  \label{eq:Lrev}
\end{equation}
It is symmetric positive semidefinite because
\begin{equation}
  v\transpose L_{\rm rev}v
  =v\transpose BCB\transpose v
  =(B\transpose v)\transpose C(B\transpose v)
  =\sum_{e=(i,j)}c_e(v_i-v_j)^2\ge 0.
  \label{eq:Lrev-psd}
\end{equation}

\begin{runningexample}[Running numerical example: reversible conductances]
For the six-node example, the stationary distribution is
\[
\begin{array}{c|cccccc}
\toprule
 i&1&2&3&4&5&6\\
\midrule
 \pi_i&0.049&0.147&0.159&0.200&0.292&0.153\\
\bottomrule
\end{array}
\]
The nonzero conductances \eqref{eq:reversible-conductance} are
\[
\begin{array}{c|cccccc}
\toprule
\{i,j\}&\{1,2\}&\{2,3\}&\{2,5\}&\{3,4\}&\{4,5\}&\{5,6\}\\
\midrule
c_{ij}&0.0491&0.0842&0.0140&0.0749&0.1250&0.1529\\
\bottomrule
\end{array}
\]
The weak edge $\{2,5\}$ is a shortcut between the left and right parts of the graph, but its conductance is much smaller than the main chain conductances.
\end{runningexample}

\subsection{A safe directed smoothness operator}

Directionality may still matter.  A safe positive semidefinite directed residual penalty is
\begin{equation}
  L_{\rm dir}:=(I-K)\transpose\Pi(I-K),
  \qquad
  \Pi:=\diag(\pi).
  \label{eq:Ldir}
\end{equation}
Then
\begin{equation}
  z\transpose L_{\rm dir}z
  =\norm{\Pi^{1/2}(I-K)z}_2^2
  =\sum_i\pi_i\left(z_i-\sum_jK_{ij}z_j\right)^2\ge0.
  \label{eq:Ldir-psd}
\end{equation}
This term says that the increment at node $i$ should be close to the directed weighted average of increments at its outgoing neighbors.  Unlike a raw directed Laplacian, \eqref{eq:Ldir} is a valid Gaussian precision contribution.

\begin{heuristicbox}[What the two graph operators do]
The reversible conductance graph answers: ``How expensive is it to move change between nodes?''  The directed residual operator answers: ``Does the increment respect the directed prediction rule encoded by $K$?''  They are complementary, not duplicates.
\end{heuristicbox}

\section{Graph OT regularization for small increments}
\label{sec:ot-regularization}

Optimal transport is naturally a geometry of nonnegative mass distributions.  Your scalar $x_i$ may or may not itself be a mass.  The safest default for arbitrary scalar attributes is therefore not to transport $x$ directly.  Instead, use OT as a \emph{geometry for increments}.  The increment $z$ is judged by the cheapest graph flux that could have produced it.

\subsection{Mobility weights}

Choose a positive reference density
\begin{equation}
  \rho^0\in\R^N_{>0},
  \qquad
  \sum_i\rho_i^0=1.
  \label{eq:reference-density}
\end{equation}
If $x^0$ is genuinely nonnegative mass, one may take $\rho^0\propto \bar x^0$.  If $x^0$ is a signed or arbitrary scalar attribute, use an exogenous density, such as uniform mass, degree-proportional mass, or precision-weighted mass.  The role of $\rho^0$ is to encode where transport is easy; it need not be the signal itself.

For edge $e=(i,j)$ define the mobility
\begin{equation}
  m_e=c_{ij}\,\theta(\rho_i^0,\rho_j^0),
  \label{eq:edge-mobility}
\end{equation}
where $\theta$ is a positive mean.  The logarithmic mean
\begin{equation}
  \theta(a,b)=\frac{a-b}{\log a-\log b},
  \qquad a\ne b,
  \qquad
  \theta(a,a)=a,
  \label{eq:log-mean}
\end{equation}
appears naturally in the Maas-Mielke discrete transport metric on finite Markov chains \cite{Maas2011,Mielke2013}.  The arithmetic mean is often sufficient in engineering implementations.

Let
\begin{equation}
  M=\diag(m_e),
  \qquad
  A_\rho=BMB\transpose.
  \label{eq:Arho}
\end{equation}
Just like a Laplacian, $A_\rho$ is symmetric positive semidefinite and annihilates constants if the graph is connected:
\[
  A_\rho\ones=0.
\]

\begin{runningexample}[Running numerical example: mobility Laplacian]
For illustration we use uniform reference density $\rho_i^0=1/6$ and the arithmetic mean.  Hence $m_e=c_e/6$.  With edge order
\[
  (1,2),(2,3),(2,5),(3,4),(4,5),(5,6),
\]
the edge mobilities are
\[
\begin{array}{c|cccccc}
\toprule
 e&(1,2)&(2,3)&(2,5)&(3,4)&(4,5)&(5,6)\\
\midrule
 m_e&0.00818&0.01403&0.00233&0.01249&0.02084&0.02548\\
\bottomrule
\end{array}
\]
The mobility Laplacian is
\[
A_\rho=\begin{bmatrix}
 0.0082&-0.0082&0&0&0&0\\
-0.0082&0.0245&-0.0140&0&-0.0023&0\\
0&-0.0140&0.0265&-0.0125&0&0\\
0&0&-0.0125&0.0333&-0.0208&0\\
0&-0.0023&0&-0.0208&0.0486&-0.0255\\
0&0&0&0&-0.0255&0.0255
\end{bmatrix}.
\]
This matrix has negative off-diagonal entries on edges and row sums equal to zero, exactly as a graph Laplacian should.
\end{runningexample}

\subsection{The conservative OT tangent norm}
\label{subsec:conservative-ot}

First suppose the increment is conservative:
\begin{equation}
  \ones\transpose z=0.
  \label{eq:conservative-increment}
\end{equation}
A flux $f\in\R^{|E|}$ assigns a signed flow to each oriented edge.  The node imbalance created by $f$ is $Bf$.  Therefore the constraint
\begin{equation}
  Bf=z
  \label{eq:Bf-z}
\end{equation}
means that the flux $f$ produces the increment $z$.

The quadratic transport energy of $f$ is
\begin{equation}
  f\transpose M^{-1}f=\sum_{e\in E}\frac{f_e^2}{m_e}.
  \label{eq:flux-energy}
\end{equation}
Large mobility $m_e$ makes edge $e$ cheap; small mobility makes it expensive.  The OT tangent norm is the least flux energy needed to produce $z$:
\begin{equation}
  \norm{z}_{\rm OT,\rho^0}^2
  :=\min_{f:\,Bf=z}f\transpose M^{-1}f.
  \label{eq:ot-primal}
\end{equation}

\paragraph{Derivation.}
Introduce a Lagrange multiplier $\phi\in\R^N$ and use the Lagrangian
\begin{equation}
  \mathcal L(f,\phi)
  =f\transpose M^{-1}f+2\phi\transpose(z-Bf).
  \label{eq:ot-lagrangian}
\end{equation}
Stationarity in $f$ gives
\begin{equation}
  \nabla_f\mathcal L=2M^{-1}f-2B\transpose\phi=0,
  \qquad\text{so}\qquad
  f=MB\transpose\phi.
  \label{eq:optimal-flux}
\end{equation}
Plugging this into the constraint gives
\begin{equation}
  BMB\transpose\phi=A_\rho\phi=z.
  \label{eq:poisson-phi}
\end{equation}
Since $A_\rho$ is singular on constants, this equation is solved on the mean-zero subspace.  The minimum-norm solution is $\phi=A_\rho\pinv z$.  The optimal value is
\begin{align}
  f\transpose M^{-1}f
  &=(MB\transpose\phi)\transpose M^{-1}(MB\transpose\phi)\notag\\
  &=\phi\transpose BMB\transpose\phi
  =\phi\transpose A_\rho\phi
  =\phi\transpose z
  =z\transpose A_\rho\pinv z.
  \label{eq:ot-tangent-value}
\end{align}
Thus
\begin{equation}
  \boxed{\norm{z}_{\rm OT,\rho^0}^2=z\transpose A_\rho\pinv z.}
  \label{eq:ot-tangent-boxed}
\end{equation}
This is the graph analogue of the Benamou-Brenier kinetic energy for infinitesimal transport \cite{Maas2011,Mielke2013,Leonard2016,PeyreCuturi2019}.

\begin{heuristicbox}[Electric-network analogy]
The equation $A_\rho\phi=z$ is a graph Poisson equation.  The increment $z$ is a vector of injections and withdrawals, $\phi$ is a potential, and $f=MB\transpose\phi$ is the induced current.  The OT tangent norm is the dissipated energy.  This is why OT regularization is global: a change at one node can be explained by withdrawals elsewhere along cheap paths.
\end{heuristicbox}

\subsection{Adding source/sink flexibility: local unbalanced OT}
\label{subsec:unbalanced-tangent}

Most scalar systems are not exactly conservative.  We therefore allow unexplained local creation/destruction $s\in\R^N$:
\begin{equation}
  z=Bf+s.
  \label{eq:z-Bf-s}
\end{equation}
Penalize source/sink activity by $\nu^{-1}\norm{s}_2^2$, where $\nu>0$ is a source/sink allowance.  Define
\begin{equation}
  \norm{z}_{\rm UOT,\rho^0,\nu}^2
  :=\min_{f,s:\,Bf+s=z}\left\{
  f\transpose M^{-1}f+\nu^{-1}\norm{s}_2^2
  \right\}.
  \label{eq:uot-primal}
\end{equation}

\paragraph{Derivation.}
Use the Lagrangian
\begin{equation}
  \mathcal L(f,s,\phi)
  =f\transpose M^{-1}f+\nu^{-1}s\transpose s
   +2\phi\transpose(z-Bf-s).
  \label{eq:uot-lagrangian}
\end{equation}
Stationarity gives
\begin{equation}
  f=MB\transpose\phi,
  \qquad
  s=\nu\phi.
  \label{eq:uot-stationarity}
\end{equation}
The constraint becomes
\begin{equation}
  (A_\rho+\nu I)\phi=z.
  \label{eq:uot-poisson}
\end{equation}
Since $\nu>0$, the matrix $A_\rho+\nu I$ is positive definite.  The optimal value is
\begin{align}
  f\transpose M^{-1}f+\nu^{-1}s\transpose s
  &=\phi\transpose A_\rho\phi+\nu\phi\transpose\phi\notag\\
  &=\phi\transpose(A_\rho+\nu I)\phi
  =\phi\transpose z
  =z\transpose(A_\rho+\nu I)^{-1}z.
  \label{eq:uot-value}
\end{align}
Therefore
\begin{equation}
  \boxed{\norm{z}_{\rm UOT,\rho^0,\nu}^2
  =z\transpose(A_\rho+\nu I)^{-1}z.}
  \label{eq:uot-boxed}
\end{equation}
Small $\nu$ makes source/sink activity expensive and approaches conservative transport.  Large $\nu$ permits more local unexplained change.

\begin{warningbox}[Meaning of $\nu$]
The parameter $\nu$ is not observation noise.  It controls how willing the prior is to explain increments by local creation/destruction instead of by graph transport.  Observation noise is already represented by $R^{-1}$.
\end{warningbox}

\section{The recommended Gaussian increment prior}
\label{sec:recommended-prior}

Combine three ingredients:
\begin{enumerate}[leftmargin=2em]
  \item local temporal precision, which says increments are usually not huge;
  \item directed graph smoothness, which says increments should respect directed local averaging;
  \item OT tangent regularization, which says increments should be explainable by low-energy graph flux plus controlled sources/sinks.
\end{enumerate}
The recommended increment precision is
\begin{equation}
  \boxed{
  Q_\Delta
  =D_\kappa+\eta L_{\rm dir}+\lambda(A_\rho+\nu I)^{-1}.}
  \label{eq:Qdelta-main}
\end{equation}
Here
\begin{equation}
  D_\kappa=\diag(\kappa_1,\ldots,\kappa_N),
  \qquad
  \kappa_i>0,
  \label{eq:Dkappa}
\end{equation}
$\eta\ge0$ controls directed smoothness, $\lambda\ge0$ controls the OT tangent penalty, and $\nu>0$ controls the source/sink allowance.  The increment prior is
\begin{equation}
  z\sim\mathcal N(0,Q_\Delta^{-1}).
  \label{eq:z-prior-final}
\end{equation}

The dense inverse in \eqref{eq:Qdelta-main} should not be formed for large $N$.  To compute $v\mapsto Q_\Delta v$, solve
\begin{equation}
  (A_\rho+\nu I)u=v
  \label{eq:A-solve}
\end{equation}
by a sparse linear solver, then return
\begin{equation}
  Q_\Delta v=D_\kappa v+\eta L_{\rm dir}v+\lambda u.
  \label{eq:Qdelta-matvec}
\end{equation}

\begin{heuristicbox}[How to read the three terms in $Q_\Delta$]
The diagonal term $D_\kappa$ says: ``large increments are unlikely locally.''  The $L_{\rm dir}$ term says: ``neighboring or directedly related increments should agree.''  The OT term says: ``positive and negative innovations should be mutually explainable by cheap graph fluxes, unless source/sink creation is allowed.''
\end{heuristicbox}

\begin{runningexample}[Running numerical example: chosen hyperparameters]
We use
\[
  \kappa_i=2.0,
  \qquad
  \eta=20.0,
  \qquad
  \lambda=0.10,
  \qquad
  \nu=0.15.
\]
The value $\eta=20$ is intentionally large in this small example to make graph propagation visible.  The resulting increment precision is approximately
\[
Q_\Delta=\begin{bmatrix}
 3.942&-1.932& 0.563& 0.000& 0.094& 0.000\\
-1.932& 7.408&-3.320& 0.916&-0.551& 0.147\\
 0.563&-3.320& 7.278&-2.957& 1.102& 0.001\\
 0.000& 0.916&-2.957& 8.331&-4.942& 1.319\\
 0.094&-0.551& 1.102&-4.942&13.003&-6.040\\
 0.000& 0.147& 0.001& 1.319&-6.040& 7.240
\end{bmatrix}.
\]
The predictive covariance is
\[
  K^-=P_0^{-1}+Q_\Delta^{-1}.
\]
Its prior standard deviations are
\[
\begin{array}{c|cccccc}
\toprule
 i&1&2&3&4&5&6\\
\midrule
 \sqrt{K^-_{ii}}&0.577&0.484&0.515&0.499&0.462&0.532\\
\bottomrule
\end{array}
\]
The columns of $K^-$ associated with the observed nodes $2$ and $5$ are
\[
K^-_{:,\{2,5\}}=
\begin{bmatrix}
 0.0828&-0.0037\\
 0.2342& 0.0032\\
 0.0865& 0.0098\\
 0.0118& 0.0808\\
 0.0032& 0.2134\\
-0.0034& 0.1215
\end{bmatrix}.
\]
These two columns already reveal the extrapolation mechanism.  A positive innovation at node $2$ should push nodes $1$ and $3$ upward because their covariances with node $2$ are positive.  A negative innovation at node $5$ should push nodes $4$ and $6$ downward because their covariances with node $5$ are positive.
\end{runningexample}

\section{Conditioning on the observed time-1 nodes}
\label{sec:conditioning}

The likelihood is
\begin{equation}
  y\given x^1\sim\mathcal N(Hx^1,R^{-1}).
  \label{eq:likelihood-final}
\end{equation}
Combining \eqref{eq:predictive-prior} and \eqref{eq:likelihood-final} gives a Gaussian posterior
\begin{equation}
  x^1\given y,\bar x^0\sim\mathcal N(\mu^+,K^+).
  \label{eq:posterior-distribution}
\end{equation}
There are two equivalent formulas.  The precision form is convenient conceptually; the covariance form is convenient when $n\ll N$.

\subsection{Precision-form derivation}

Let $Q^-=(K^-)^{-1}$.  Up to an additive constant independent of $x$, the negative log posterior is
\begin{align}
  \Phi(x)
  &=\half(x-\mu^-)\transpose Q^-(x-\mu^-)
    +\half(y-Hx)\transpose R(y-Hx)\notag\\
  &=\half x\transpose(Q^-+H\transpose R H)x
    -x\transpose(Q^-\mu^-+H\transpose R y)
    +\text{constant}.
  \label{eq:posterior-quadratic-expanded}
\end{align}
Completing the square gives
\begin{equation}
  \boxed{Q^+=Q^-+H\transpose RH,}
  \label{eq:Qplus}
\end{equation}
\begin{equation}
  \boxed{\mu^+=(Q^+)^{-1}\left(Q^-\mu^-+H\transpose R y\right),}
  \label{eq:muplus-precision}
\end{equation}
\begin{equation}
  \boxed{K^+=(Q^+)^{-1}.}
  \label{eq:Kplus-precision}
\end{equation}
This is the direct Bayesian update formula.

\subsection{Covariance-form derivation}

When $n\ll N$, invert only an $n\times n$ matrix.  Define the innovation
\begin{equation}
  d:=y-H\mu^-.
  \label{eq:innovation-d}
\end{equation}
The covariance of $y$ before seeing the observation is
\begin{equation}
  S_y:=\Cov(y\given \bar x^0)=HK^-H\transpose+R^{-1}.
  \label{eq:innovation-covariance-final}
\end{equation}
The joint prior distribution of $x^1$ and $y$ is Gaussian with
\begin{equation}
  \E\begin{bmatrix}x^1\\y\end{bmatrix}
  =\begin{bmatrix}\mu^-\\H\mu^-\end{bmatrix},
  \qquad
  \Cov\begin{bmatrix}x^1\\y\end{bmatrix}
  =\begin{bmatrix}
    K^-&K^-H\transpose\\
    HK^-&S_y
  \end{bmatrix}.
  \label{eq:joint-x-y}
\end{equation}
The standard conditioning formula for a joint Gaussian gives
\begin{equation}
  \boxed{
  \mu^+=\mu^-+K^-H\transpose S_y^{-1}(y-H\mu^-),}
  \label{eq:muplus-covariance}
\end{equation}
\begin{equation}
  \boxed{
  K^+=K^- - K^-H\transpose S_y^{-1}HK^-.}
  \label{eq:Kplus-covariance}
\end{equation}
The same identities follow from \eqref{eq:Qplus} using the Woodbury matrix identity.

\paragraph{Interpretation.}  The vector
\begin{equation}
  \alpha:=S_y^{-1}(y-H\mu^-)
  \label{eq:alpha-definition}
\end{equation}
contains the precision-weighted observed innovations.  The product $K^-H\transpose\alpha$ spreads these innovations to all nodes using graph-informed prior covariances.

\begin{runningexample}[Running numerical example: the update]
For the chosen hyperparameters,
\[
  S_y=
  \begin{bmatrix}
    0.2442&0.0032\\
    0.0032&0.2290
  \end{bmatrix},
  \qquad
  d=\begin{bmatrix}0.35\\-0.25\end{bmatrix}.
\]
Thus
\[
  \alpha=S_y^{-1}d
  =\begin{bmatrix}1.4478\\-1.1118\end{bmatrix}.
\]
The posterior mean is
\[
  \mu^+=\bar x^0+K^-_{:,\{2,5\}}\alpha.
\]
For instance,
\begin{align*}
  \mu_3^+
  &=1.95+0.0865(1.4478)+0.0098(-1.1118)\\
  &\approx 2.064,\\[1mm]
  \mu_6^+
  &=1.10+(-0.0034)(1.4478)+0.1215(-1.1118)\\
  &\approx 0.960.
\end{align*}
Node $3$ moves upward because it is close to the positive innovation at node $2$.  Node $6$ moves downward because it is close to the negative innovation at node $5$.
\end{runningexample}

\subsection{Block formula for observed and unobserved nodes}
\label{subsec:block-formula}

Partition the nodes into observed nodes $S$ and unobserved nodes $U$.  For notational simplicity, assume $H$ selects the observed block directly.  Partition
\[
  x=\begin{bmatrix}x_S\\x_U\end{bmatrix},
  \qquad
  \mu^-=\begin{bmatrix}\mu_S^-\\\mu_U^-\end{bmatrix},
  \qquad
  K^-=
  \begin{bmatrix}
    K^-_{SS}&K^-_{SU}\\
    K^-_{US}&K^-_{UU}
  \end{bmatrix}.
\]
Then \eqref{eq:muplus-covariance} and \eqref{eq:Kplus-covariance} imply
\begin{equation}
  \boxed{
  \mu_U^+=\mu_U^-+K^-_{US}(K^-_{SS}+R^{-1})^{-1}(y-\mu_S^-),}
  \label{eq:block-mean-final}
\end{equation}
\begin{equation}
  \boxed{
  K^+_{UU}=K^-_{UU}-K^-_{US}(K^-_{SS}+R^{-1})^{-1}K^-_{SU}.}
  \label{eq:block-cov-final}
\end{equation}
Thus unobserved nodes are updated only through their prior covariance with observed nodes.  If node $i$ is weakly coupled to all observed nodes, the row $K^-_{iS}$ is small; its posterior mean remains close to its baseline prediction and its uncertainty barely shrinks.

\subsection{Posterior marginal precision}

Once $K^+$ is available, the reported precision is
\begin{equation}
  \boxed{\pi_i^+=\frac{1}{K^+_{ii}}.}
  \label{eq:marginal-precision-final}
\end{equation}
The corresponding approximate $95\%$ posterior credible interval is
\begin{equation}
  \mu_i^+\pm1.96\sqrt{K^+_{ii}}.
  \label{eq:credible-interval}
\end{equation}

\begin{runningexample}[Running numerical example: final posterior table]
The posterior values and marginal precisions are
\[
\begin{array}{c|ccccc}
\toprule
 i&\bar x_i^0&\text{time-1 obs.}&\mu_i^+&\sqrt{K^+_{ii}}&\pi_i^+=1/K^+_{ii}\\
\midrule
1&2.00&-&2.124&0.552&  3.286\\
2&2.10&2.45&2.436&0.098&104.271\\
3&1.95&-&2.064&0.484&  4.274\\
4&1.05&-&0.977&0.468&  4.556\\
5&0.90&0.65&0.667&0.121& 68.687\\
6&1.10&-&0.960&0.468&  4.568\\
\bottomrule
\end{array}
\]
The observed nodes $2$ and $5$ have high posterior precision because their observation precisions are high.  The unobserved nodes move in the graph-plausible direction but remain much less precise.  For example, node $3$ moves from $1.95$ to $2.064$, but its posterior standard deviation is still about $0.484$.

For comparison, the diagonal entries of the full posterior precision matrix are
\[
\begin{array}{c|cccccc}
\toprule
 i&1&2&3&4&5&6\\
\midrule
 Q^+_{ii}&3.314&105.348&4.666&4.985&70.995&4.689\\
 \pi_i^+&3.286&104.271&4.274&4.556&68.687&4.568\\
\bottomrule
\end{array}
\]
They are close in this small example, but they are not the same quantities.  In stronger-correlation regimes the difference can be large.
\end{runningexample}

\section{How the OT strength changes extrapolation}
\label{sec:lambda-effects}

The OT weight $\lambda$ controls how strongly the prior penalizes increments that require high transport/source energy.  Holding the other hyperparameters fixed in the running example gives:
\[
\begin{array}{c|cccccc|cccccc}
\toprule
&\multicolumn{6}{c|}{\text{posterior means}}&\multicolumn{6}{c}{\text{posterior standard deviations}}\\
\lambda&1&2&3&4&5&6&1&2&3&4&5&6\\
\midrule
0.00&2.154&2.438&2.078&0.972&0.665&0.937&0.599&0.098&0.508&0.489&0.121&0.489\\
0.05&2.138&2.437&2.071&0.975&0.666&0.949&0.573&0.098&0.495&0.478&0.121&0.478\\
0.10&2.124&2.436&2.064&0.977&0.667&0.960&0.552&0.098&0.484&0.468&0.121&0.468\\
0.50&2.065&2.430&2.028&0.993&0.674&1.014&0.445&0.097&0.424&0.416&0.119&0.410\\
1.00&2.036&2.424&2.003&1.008&0.679&1.047&0.381&0.096&0.385&0.380&0.118&0.370\\
\bottomrule
\end{array}
\]
As $\lambda$ increases, the prior admits fewer high-energy increment patterns.  The unobserved posterior means are pulled closer to the baseline and the posterior standard deviations shrink.  The observed nodes remain close to the data because their observation precisions are high.

\begin{heuristicbox}[Choosing $\lambda$]
A very small $\lambda$ lets directed smoothness dominate, so observed innovations can propagate widely.  A very large $\lambda$ can make the model conservative and stiff: it suppresses many increment patterns, including some that may be real.  In practice, tune $\lambda$ by marginal likelihood or held-out calibration, not by visual appeal alone.
\end{heuristicbox}

\section{What to compute on a large sparse graph}
\label{sec:large-sparse-computation}

For large $N$, avoid dense $N\times N$ inverses.  The key observation is that $n\ll N$, so the posterior mean needs only the $n$ columns of $K^-$ corresponding to observed nodes.

\subsection{Computing the posterior mean}

The predictive covariance is
\[
  K^-=P_0^{-1}+Q_\Delta^{-1}.
\]
To form $K^-H\transpose$, solve
\begin{equation}
  Q_\Delta u_a=H\transpose e_a,
  \qquad a=1,\ldots,n.
  \label{eq:n-solves}
\end{equation}
Then
\begin{equation}
  K^-H\transpose
  =P_0^{-1}H\transpose+[u_1\ \cdots\ u_n].
  \label{eq:KminusH-columns}
\end{equation}
After this, form the small matrix
\begin{equation}
  S_y=HK^-H\transpose+R^{-1}
  \label{eq:small-Sy}
\end{equation}
and solve
\begin{equation}
  S_y\alpha=y-H\mu^-.
  \label{eq:small-alpha}
\end{equation}
Finally,
\begin{equation}
  \mu^+=\mu^-+(K^-H\transpose)\alpha.
  \label{eq:large-mu}
\end{equation}

\subsection{Computing posterior variances}

The diagonal of \eqref{eq:Kplus-covariance} is
\begin{equation}
  K^+_{ii}=K^-_{ii}-k_i\transpose S_y^{-1}k_i,
  \qquad
  k_i:=HK^-e_i\in\R^n.
  \label{eq:posterior-diag-formula}
\end{equation}
The vector $k_i$ is the $i$th row of $K^-H\transpose$, transposed.  Thus the update term is cheap once $K^-H\transpose$ and $S_y^{-1}$ are available.  The remaining challenge is the prior diagonal $K^-_{ii}$:
\begin{equation}
  K^-_{ii}=1/p_i^0+(Q_\Delta^{-1})_{ii}.
  \label{eq:prior-diag}
\end{equation}
Common choices are:
\[
\begin{array}{p{0.30\linewidth}p{0.60\linewidth}}
\toprule
\text{Method}&\text{When it is useful}\\
\midrule
Sparse Cholesky selected inverse&Moderate $N$ and acceptable fill-in.  Gives accurate selected inverse diagonals.\\
Hutchinson probing&Matrix-free estimate of $\diag(Q_\Delta^{-1})$ using random signs $v_m$ and solves $Q_\Delta u_m=v_m$.\\
Stochastic Lanczos quadrature&Useful when many uncertainty summaries are required and matrix-vector products are cheap.\\
Exact subset solves&If only some nodes need uncertainties, solve $Q_\Delta u=e_i$ for those nodes.\\
\bottomrule
\end{array}
\]

\paragraph{Hutchinson diagonal estimator.}  If $v_m$ are independent Rademacher vectors, meaning each entry is $+1$ or $-1$ with probability $1/2$, then
\begin{equation}
  \E[v_m\odot(Q_\Delta^{-1}v_m)]=\diag(Q_\Delta^{-1}),
  \label{eq:hutchinson-unbiased}
\end{equation}
where $\odot$ is entrywise multiplication.  Hence
\begin{equation}
  \widehat{\diag(Q_\Delta^{-1})}
  =\frac1M\sum_{m=1}^M v_m\odot u_m,
  \qquad
  Q_\Delta u_m=v_m,
  \label{eq:hutchinson-estimator}
\end{equation}
is unbiased.  Its variance falls like $1/M$.

\section{Hyperparameter calibration}
\label{sec:calibration}

The main hyperparameters are
\begin{equation}
  \theta=(\kappa,\eta,\lambda,\nu,\rho^0,\text{graph normalization choices}).
  \label{eq:theta}
\end{equation}
They should be calibrated whenever possible.  Two practical methods are empirical Bayes and held-out predictive calibration.

\subsection{Empirical Bayes by marginal likelihood}

Under the Gaussian model,
\begin{equation}
  y\given\bar x^0,\theta
  \sim\mathcal N(H\mu^-_\theta,S_{y,\theta}),
  \qquad
  S_{y,\theta}=HK^-_\theta H\transpose+R^{-1}.
  \label{eq:marginal-y}
\end{equation}
The log marginal likelihood is
\begin{equation}
  \ell(\theta)
  =-\half d_\theta\transpose S_{y,\theta}^{-1}d_\theta
   -\half\log\det S_{y,\theta}
   -\frac n2\log(2\pi),
  \qquad
  d_\theta=y-H\mu^-_\theta.
  \label{eq:log-marginal}
\end{equation}
Since $S_{y,\theta}$ is only $n\times n$, evaluating \eqref{eq:log-marginal} is often feasible even when $N$ is large, as long as the required covariance columns are available.

\paragraph{Gradient derivation.}  Suppose $\mu^-_\theta$ does not depend on a scalar hyperparameter $\theta_j$; this holds for $\kappa,\eta,\lambda,\nu$ if the drift is fixed.  Let
\[
  S=S_{y,\theta},
  \qquad
  d=d_\theta,
  \qquad
  \alpha=S^{-1}d.
\]
Using
\[
  \partial_\theta S^{-1}=-S^{-1}(\partial_\theta S)S^{-1},
  \qquad
  \partial_\theta\log\det S=\tr(S^{-1}\partial_\theta S),
\]
we get
\begin{align}
  \partial_{\theta_j}\ell
  &=\half d\transpose S^{-1}(\partial_{\theta_j}S)S^{-1}d
    -\half\tr(S^{-1}\partial_{\theta_j}S)\notag\\
  &=\half\alpha\transpose(\partial_{\theta_j}S)\alpha
    -\half\tr(S^{-1}\partial_{\theta_j}S).
  \label{eq:marginal-gradient}
\end{align}
If $\mu^-_\theta$ also depends on $\theta_j$, then add
\begin{equation}
  (\partial_{\theta_j}H\mu^-_\theta)\transpose S^{-1}d
  \label{eq:marginal-gradient-mean-term}
\end{equation}
with the appropriate sign after differentiating $d=y-H\mu^-_\theta$.

\subsection{Held-out calibration}

If the observed subset is large enough, split it into a fitting set and a validation set.  For a held-out observed node $i$, the posterior predictive distribution is
\begin{equation}
  y_i\given y_{\rm fit}
  \sim\mathcal N(\mu_i^+,K^+_{ii}+r_i^{-1}).
  \label{eq:heldout-predictive}
\end{equation}
Define the standardized residual
\begin{equation}
  \zeta_i=\frac{y_i-\mu_i^+}{\sqrt{K^+_{ii}+r_i^{-1}}}.
  \label{eq:standardized-residual-final}
\end{equation}
If the model is calibrated, the $\zeta_i$ should look roughly like standard normal draws.  If $|\zeta_i|$ is systematically too large, the posterior is overconfident or biased.  If $|\zeta_i|$ is systematically too small, the posterior is underconfident.

\section{Full nonlinear unbalanced OT alternative}
\label{sec:full-nonlinear-uot}

The linearized model above is the recommended default for small changes.  It is especially appropriate when $x_i$ is an arbitrary scalar attribute rather than a mass.  If the scalar values are genuinely nonnegative masses, or can be faithfully mapped to nonnegative masses, one can instead use a full unbalanced OT penalty.

Let $T:\R\to\R_+$ be a positivity map applied coordinatewise.  Examples are
\begin{equation}
  T(x_i)=x_i\quad\text{if }x_i\ge0,
  \qquad
  T(x_i)=\log(1+\exp((x_i-a)/s))+\delta,
  \label{eq:positivity-map}
\end{equation}
or a two-channel signed representation
\begin{equation}
  T_+(x_i)=[x_i]_+,
  \qquad
  T_-(x_i)=[-x_i]_+.
  \label{eq:two-channel}
\end{equation}
Let
\[
  a=T(\bar x^0),
  \qquad
  b=T(x).
\]
With a graph shortest-path cost matrix $C_{ij}$ and entropic parameter $\eps>0$, one form of unbalanced OT is
\begin{align}
  \operatorname{UOT}_{C,\eps,\rho_a,\rho_b}(a,b)
  =\min_{\Gamma\ge0}\;&
  \inner{C}{\Gamma}
  +\eps\sum_{ij}\Gamma_{ij}\left(\log\frac{\Gamma_{ij}}{a_ib_j}-1\right)\notag\\
  &+\rho_a\,\KL(\Gamma\ones\,\|\,a)
  +\rho_b\,\KL(\Gamma\transpose\ones\,\|\,b).
  \label{eq:full-uot}
\end{align}
Here $\Gamma_{ij}$ is the amount of mass transported from $i$ to $j$.  The two KL terms allow the transported row and column masses to differ from $a$ and $b$, hence the transport is unbalanced \cite{Chizat2018,Sejourne2022}.

The nonlinear MAP estimator is
\begin{align}
  \widehat x_{\rm MAP}
  =\argmin_x\Bigl\{&
  \half\norm{Hx-y}_R^2
  +\half\norm{x-\bar x^0-m_\Delta}_{D_\kappa}^2\notag\\
  &+\frac{\eta}{2}(x-\bar x^0-m_\Delta)\transpose
    L_{\rm dir}(x-\bar x^0-m_\Delta)\notag\\
  &+\lambda\operatorname{UOT}_{C,\eps,\rho_a,\rho_b}
    (T(\bar x^0),T(x))\Bigr\}.
  \label{eq:nonlinear-uot-map}
\end{align}
This is generally non-Gaussian.  A standard uncertainty approximation is the Laplace approximation.  Let
\begin{equation}
  \Phi(x)=-\log p(x\given y,\bar x^0)
  \label{eq:Phi-nonlinear}
\end{equation}
be the objective in \eqref{eq:nonlinear-uot-map}.  At the MAP,
\begin{equation}
  Q_{\rm Lap}=\nabla^2\Phi(\widehat x_{\rm MAP}),
  \qquad
  K_{\rm Lap}=Q_{\rm Lap}^{-1},
  \qquad
  \pi_i^{\rm Lap}=1/(K_{\rm Lap})_{ii}.
  \label{eq:Laplace}
\end{equation}
Entropic OT is differentiable and can be solved with Sinkhorn-type algorithms \cite{Cuturi2013,Chizat2018,Flamary2021,Cuturi2022OTT}.  However, posterior uncertainty can be skewed when observations are very sparse or when $T$ is strongly nonlinear, so Laplace precision should be validated by predictive calibration.

\begin{warningbox}[When not to use full OT directly]
If $x_i$ is a rating, temperature, score, residual, signed coefficient, or other arbitrary scalar, then treating $x$ as a mass distribution is usually conceptually wrong.  Use the tangent OT regularizer on the increment with a separate positive mobility density $\rho^0$, or split signed quantities into positive and negative channels.
\end{warningbox}

\section{Recommended workflow}
\label{sec:workflow}

\begin{enumerate}[leftmargin=2em]
  \item \textbf{Normalize and inspect the graph.}  Build $K$ from $W$, handle isolated nodes, and identify connected components.
  \item \textbf{Construct valid symmetric operators.}  Compute $\pi$, conductances $c_{ij}$, $A_\rho$, and optionally $L_{\rm dir}$.
  \item \textbf{Choose the mobility density.}  Use $\rho^0\propto\bar x^0$ only for genuine nonnegative mass.  Otherwise use a positive exogenous density.
  \item \textbf{Build the increment precision.}  Start with \eqref{eq:Qdelta-main}.  For a first model, take $D_\kappa=\kappa I$ and tune $(\kappa,\eta,\lambda,\nu)$.
  \item \textbf{Propagate baseline uncertainty.}  Set $K^-=P_0^{-1}+Q_\Delta^{-1}$ and $\mu^-=\bar x^0+m_\Delta$.
  \item \textbf{Condition on time-$1$ observations.}  Use \eqref{eq:muplus-covariance} and \eqref{eq:Kplus-covariance}, exploiting $n\ll N$.
  \item \textbf{Report means and marginal precisions.}  Output $\mu_i^+$ and $1/K^+_{ii}$, not just a MAP value.
  \item \textbf{Calibrate.}  Use marginal likelihood, held-out observations, and standardized residuals to check whether uncertainty is credible.
\end{enumerate}

\section{Failure modes and safeguards}
\label{sec:failure-modes}

\begin{enumerate}[label=(\alph*),leftmargin=2em]
  \item \textbf{Regularizing the absolute field.}  This can erase real baseline structure.  Regularize $x^1-x^0$ unless you truly believe the absolute signal is graph-smooth.
  \item \textbf{Using a nonsymmetric directed Laplacian as a precision.}  Gaussian precisions must be symmetric positive semidefinite.  Use $L_{\rm rev}$, $L_{\rm dir}$, or another symmetric construction.
  \item \textbf{Using OT directly on signed scalars.}  OT moves nonnegative mass.  For arbitrary scalars, use OT as a geometry for increments or use a positive transform with care.
  \item \textbf{Ignoring baseline precision.}  Finite $p_i^0$ contributes uncertainty at time $1$.  Add $P_0^{-1}$ to the predictive covariance.
  \item \textbf{Reporting conditional precision as marginal precision.}  Report $1/K^+_{ii}$.
  \item \textbf{Disconnected components with no observations.}  Such components cannot learn from time-$1$ data.  Their posterior should remain close to baseline with uncertainty governed by $p_i^0$ and the increment prior.
  \item \textbf{Over-tuning to the observed subset.}  If $n$ is small, empirical Bayes may prefer overconfident hyperparameters.  Use conservative priors and held-out calibration when possible.
\end{enumerate}

\appendix

\section{Gaussian conditioning from first principles}
\label{app:gaussian-conditioning}

Let
\[
  \begin{bmatrix}x\\y\end{bmatrix}
  \sim\mathcal N\left(
  \begin{bmatrix}\mu_x\\\mu_y\end{bmatrix},
  \begin{bmatrix}
    K_{xx}&K_{xy}\\
    K_{yx}&K_{yy}
  \end{bmatrix}
  \right).
\]
Assume $K_{yy}$ is invertible.  We derive
\begin{equation}
  x\given y
  \sim\mathcal N\left(
  \mu_x+K_{xy}K_{yy}^{-1}(y-\mu_y),
  K_{xx}-K_{xy}K_{yy}^{-1}K_{yx}
  \right).
  \label{eq:joint-conditioning-formula}
\end{equation}

Set
\begin{equation}
  \tilde x=x-\mu_x-K_{xy}K_{yy}^{-1}(y-\mu_y).
  \label{eq:tilde-x}
\end{equation}
Then $\E[\tilde x]=0$ and
\begin{align}
  \Cov(\tilde x,y)
  &=\Cov(x,y)-K_{xy}K_{yy}^{-1}\Cov(y,y)
  =K_{xy}-K_{xy}=0.
  \label{eq:tilde-uncorrelated-y}
\end{align}
Because $(\tilde x,y)$ is jointly Gaussian, zero covariance implies independence.  Hence conditioning on $y$ does not change the distribution of $\tilde x$.  Its covariance is
\begin{align}
  \Cov(\tilde x)
  &=K_{xx}-K_{xy}K_{yy}^{-1}K_{yx}.
  \label{eq:tilde-covariance}
\end{align}
Rearranging \eqref{eq:tilde-x} gives \eqref{eq:joint-conditioning-formula}.  Applying this formula with
\[
  K_{xx}=K^- ,
  \quad
  K_{xy}=K^-H\transpose,
  \quad
  K_{yy}=HK^-H\transpose+R^{-1}
\]
produces \eqref{eq:muplus-covariance} and \eqref{eq:Kplus-covariance}.

\section{Woodbury identity connection}
\label{app:woodbury}

The Woodbury identity says
\begin{equation}
  (A+UCV)^{-1}
  =A^{-1}-A^{-1}U(C^{-1}+VA^{-1}U)^{-1}VA^{-1}.
  \label{eq:woodbury}
\end{equation}
Choose
\[
  A=Q^-=(K^-)^{-1},
  \qquad
  U=H\transpose,
  \qquad
  C=R,
  \qquad
  V=H.
\]
Then
\begin{align}
  (Q^-+H\transpose RH)^{-1}
  &=K^- - K^-H\transpose(R^{-1}+HK^-H\transpose)^{-1}HK^-.
  \label{eq:woodbury-posterior}
\end{align}
This is exactly \eqref{eq:Kplus-covariance}.  Multiplying the precision-form mean \eqref{eq:muplus-precision} by the same identity gives \eqref{eq:muplus-covariance}.

\section{Why marginal and conditional precision differ}
\label{app:precision-difference}

Consider a two-dimensional posterior covariance
\begin{equation}
  K=\begin{bmatrix}
    \sigma_1^2&\rho\sigma_1\sigma_2\\
    \rho\sigma_1\sigma_2&\sigma_2^2
  \end{bmatrix},
  \qquad |\rho|<1.
  \label{eq:two-var-cov}
\end{equation}
The marginal precision of $x_1$ is
\begin{equation}
  \pi_1^{\rm marg}=1/\sigma_1^2.
  \label{eq:marginal-two}
\end{equation}
But
\begin{equation}
  Q=K^{-1}
  =\frac{1}{\sigma_1^2\sigma_2^2(1-\rho^2)}
  \begin{bmatrix}
    \sigma_2^2&-\rho\sigma_1\sigma_2\\
    -\rho\sigma_1\sigma_2&\sigma_1^2
  \end{bmatrix},
  \label{eq:two-var-precision}
\end{equation}
so
\begin{equation}
  Q_{11}=\frac{1}{\sigma_1^2(1-\rho^2)}.
  \label{eq:conditional-two}
\end{equation}
Since $1/(1-\rho^2)>1$, the conditional precision $Q_{11}$ is larger than the marginal precision whenever $\rho\ne0$.  This is why reporting $Q^+_{ii}$ can overstate certainty.

\section{Dual form of the OT tangent norm}
\label{app:ot-dual}

The conservative primal problem is
\[
  \min_{Bf=z} f\transpose M^{-1}f.
\]
Using the stationarity $f=MB\transpose\phi$, the optimal flux is a gradient flow over the graph.  An equivalent dual form is obtained by substituting the minimized $f$ into the Lagrangian:
\begin{align}
  \inf_f\mathcal L(f,\phi)
  &=\inf_f\{f\transpose M^{-1}f-2\phi\transpose Bf+2\phi\transpose z\}\notag\\
  &=2\phi\transpose z-\phi\transpose BMB\transpose\phi.
  \label{eq:ot-dual-derivation}
\end{align}
Therefore
\begin{equation}
  \norm{z}_{\rm OT,\rho^0}^2
  =\sup_{\phi}\left\{2\phi\transpose z-\phi\transpose A_\rho\phi\right\}.
  \label{eq:ot-dual}
\end{equation}
The maximizer satisfies $A_\rho\phi=z$, and the maximum equals $z\transpose A_\rho\pinv z$.  The unbalanced version similarly gives
\begin{equation}
  \norm{z}_{\rm UOT,\rho^0,\nu}^2
  =\sup_{\phi}\left\{2\phi\transpose z-\phi\transpose(A_\rho+\nu I)\phi\right\}.
  \label{eq:uot-dual}
\end{equation}

\section{Minimal pseudocode}
\label{app:pseudocode}

\begin{verbatim}
Input: directed weights W, baseline x0_bar, baseline precisions p0,
       observed nodes S, observations y, observation precisions r

1. K = row_normalize(W)
2. pi = stationary_distribution(K)
3. c_ij = 0.5 * (pi_i K_ij + pi_j K_ji)
4. Build an incidence matrix B on the undirected support of c_ij
5. Choose positive reference density rho0
6. m_e = c_e * mean(rho0_i, rho0_j)
7. A = B diag(m_e) B^T
8. L_dir = (I-K)^T diag(pi) (I-K)
9. Define Q_delta(v):
       u = solve(A + nu I, v)
       return D_kappa v + eta * L_dir v + lambda * u
10. For each observed node s_a:
       h_a = basis vector at node s_a
       u_a = solve(Q_delta, h_a)
       column_a = P0^{-1} h_a + u_a
11. Let Cobs = [column_1 ... column_n] = K_minus H^T
12. S_y = H Cobs + R^{-1}
13. alpha = solve(S_y, y - H mu_minus)
14. mu_plus = mu_minus + Cobs alpha
15. Estimate diag(K_minus), then for each node i:
       k_i = row i of Cobs, transposed
       var_plus_i = K_minus_ii - k_i^T solve(S_y, k_i)
       precision_plus_i = 1 / var_plus_i

Output: posterior means mu_plus and marginal precisions precision_plus_i
\end{verbatim}

\begin{thebibliography}{99}

\bibitem{RasmussenWilliams2006}
C. E. Rasmussen and C. K. I. Williams.
\newblock \emph{Gaussian Processes for Machine Learning}.
\newblock MIT Press, 2006.

\bibitem{RueHeld2005}
H. Rue and L. Held.
\newblock \emph{Gaussian Markov Random Fields: Theory and Applications}.
\newblock Chapman and Hall/CRC, 2005.

\bibitem{Zhang2015GSP}
C. Zhang, D. Florencio, and P. A. Chou.
\newblock Graph signal processing: A probabilistic framework.
\newblock Microsoft Research Technical Report MSR-TR-2015-31, 2015.

\bibitem{Ng2018GraphGP}
Y. C. Ng, N. Colombo, and R. Silva.
\newblock Bayesian semi-supervised learning with graph Gaussian processes.
\newblock \emph{Advances in Neural Information Processing Systems}, 2018.

\bibitem{RompelbergSchaub2025}
L. Rompelberg and M. T. Schaub.
\newblock A Bayesian perspective on uncertainty quantification for estimated graph signals.
\newblock arXiv:2502.12971, 2025.

\bibitem{Dong2016LaplacianLearning}
X. Dong, D. Thanou, P. Frossard, and P. Vandergheynst.
\newblock Learning Laplacian matrix in smooth graph signal representations.
\newblock \emph{IEEE Transactions on Signal Processing}, 64(23):6160--6173, 2016.

\bibitem{Maas2011}
J. Maas.
\newblock Gradient flows of the entropy for finite Markov chains.
\newblock \emph{Journal of Functional Analysis}, 261(8):2250--2292, 2011.

\bibitem{Mielke2013}
A. Mielke.
\newblock Geodesic convexity of the relative entropy in reversible Markov chains.
\newblock \emph{Calculus of Variations and Partial Differential Equations}, 48:1--31, 2013.

\bibitem{Leonard2016}
C. Leonard.
\newblock Lazy random walks and optimal transport on graphs.
\newblock \emph{Annals of Probability}, 44(3):1864--1915, 2016.

\bibitem{Erbar2020}
M. Erbar, M. Fathi, V. Laschos, and A. Schlichting.
\newblock Gradient flow structure for McKean-Vlasov equations on discrete spaces.
\newblock \emph{Discrete and Continuous Dynamical Systems}, 40(12):6799--6839, 2020.

\bibitem{Cuturi2013}
M. Cuturi.
\newblock Sinkhorn distances: Lightspeed computation of optimal transport.
\newblock \emph{Advances in Neural Information Processing Systems}, 2013.

\bibitem{PeyreCuturi2019}
G. Peyre and M. Cuturi.
\newblock Computational optimal transport.
\newblock \emph{Foundations and Trends in Machine Learning}, 11(5--6):355--607, 2019.

\bibitem{Chizat2018}
L. Chizat, G. Peyre, B. Schmitzer, and F.-X. Vialard.
\newblock Scaling algorithms for unbalanced optimal transport problems.
\newblock \emph{Mathematics of Computation}, 87(314):2563--2609, 2018.

\bibitem{Sejourne2022}
T. Sejourne, F.-X. Vialard, and G. Peyre.
\newblock Unbalanced optimal transport, from theory to numerics.
\newblock arXiv:2211.08775, 2022.

\bibitem{Flamary2014LOT}
R. Flamary, N. Courty, D. Tuia, and A. Rakotomamonjy.
\newblock Optimal transport with Laplacian regularization.
\newblock NIPS Workshop on Optimal Transport and Machine Learning, 2014.

\bibitem{Flamary2021}
R. Flamary, N. Courty, A. Gramfort, M. Alaya, A. Boisbunon, S. Chambon, L. Chapel, A. Corenflos, K. Fatras, N. Fournier, L. Gautheron, N. T. H. Gayraud, H. Janati, A. Rakotomamonjy, I. Redko, A. Rolet, A. Schutz, V. Seguy, D. J. Sutherland, R. Tavenard, A. Tong, and T. Vayer.
\newblock POT: Python Optimal Transport.
\newblock \emph{Journal of Machine Learning Research}, 22(78):1--8, 2021.

\bibitem{Cuturi2022OTT}
M. Cuturi, L. Meng-Papaxanthos, Y. Tian, C. Bunne, G. Davis, and O. Teboul.
\newblock Optimal Transport Tools (OTT): A JAX toolbox for all things Wasserstein.
\newblock arXiv:2201.12324, 2022.

\bibitem{KarlssonRingh2018}
J. Karlsson and A. Ringh.
\newblock Optimal mass transport for regularizing inverse problems using generalized Sinkhorn iterations.
\newblock In \emph{Proceedings of MTNS}, 2018.

\bibitem{BrediesFanzon2020}
K. Bredies and S. Fanzon.
\newblock An optimal transport approach for solving dynamic inverse problems in spaces of measures.
\newblock \emph{ESAIM: Mathematical Modelling and Numerical Analysis}, 54(6):2351--2382, 2020.

\end{thebibliography}

\end{document}
